\section{Compilation and Retrieval}
\label{sec:method}

\subsection{A shared scan for both certificates}

The compiler receives claims, date bounds, and accepted directed replacement pairs.
It initializes both boundaries to infinity.
Each accepted pair updates at most two minima.
Figure~\ref{fig:compiler} gives the complete boundary update.
The same witness can supply both boundaries.
If several witnesses tie, their identifiers determine a stable choice.

\begin{figure}[t]
\centering
\fbox{\begin{minipage}{.92\columnwidth}
\small
\textbf{Input:} claims and accepted pairs $(d,c)$.\\
\textbf{Output:} $(p_c,r_c)$ and their source witnesses.
\begin{enumerate}
\item Set $p_c=r_c=\infty$ for every claim.
\item For each accepted pair with $d\ne c$:
\item\hspace{.4em} If $\ell_d>u_c$, minimize $p_c$ using $u_d$.
\item\hspace{.4em} If $u_d>\ell_c$, minimize $r_c$ using $\ell_d$.
\item Store the supplying witness for each minimum.
\item Break boundary ties by source identifier.
\end{enumerate}
Steps 3--4 execute inside the pair scan.
\end{minipage}}
\caption{Two-boundary compilation. Every accepted pair is examined once. The stored minima answer every query window.}
\label{fig:compiler}
\end{figure}

Let $n$ count claims and $m$ count supplied accepted pairs.
Boundary compilation takes $O(n+m)$ time and stores $O(n)$ records.
A decision for one claim takes $O(1)$ time.
A complete support mask takes $O(n)$ time.
The reference gate interface can discover pairs through a complete directed scan.
This scan makes $n(n-1)$ gate calls.
The certificate representation does not require storing those calls or their rejected pairs.

Candidate discovery and boundary compilation serve different purposes.
An inverted index can reduce gate calls when its postings cover every accepted pair.
Incomplete blocking can omit a replacement and change a boundary.
Our implementation provides a complete scan as its executable specification.
The pair-stream interface supports an existing extractor or a complete index.
Both interfaces return the same certificate format.

\subsection{Filtering a fixed ranking}

Let $\pi_q$ be a complete relevance ranking for question $q$.
Let $M(q)$ contain claims with the chosen support status.
Retrieval returns
\begin{equation}
 R(q)=\operatorname{Top}_k\langle c\in\pi_q:c\in M(q)\rangle.
\label{eq:context}
\end{equation}
The filter preserves the relative order of retained claims.
It runs before the context budget.
Possible-support retrieval uses the union of guaranteed and unresolved claims.
Guaranteed-support retrieval uses only guaranteed claims.
The compiler supplies both masks without changing relevance scores.

These policies serve different evidence needs.
Possible support is the smallest pool that preserves every admissible supporting claim.
Its complement identifies evidence that can be excluded across all permitted timelines.
Guaranteed support contains claims whose temporal support survives every permitted date assignment.
Unresolved claims retain their source text and bounds for downstream reasoning.
Neither policy turns temporal support into a relevance score.


\subsection{Certifying an unchanged context}
A reader uses a ranked context with a fixed budget $k$.
The complete ranking fixes all score ties.
Let $P(Q)$ and $H(Q)$ denote possible and guaranteed support.
Write $\operatorname{Top}_k(S)$ for the first $k$ ranked claims in $S$.

\begin{theorem}[Exact context stability]
\label{thm:context-stability}
Every allowed date assignment returns the same context exactly when
\begin{equation}
 \operatorname{Top}_k(P(Q))=\operatorname{Top}_k(H(Q)).
\label{eq:context-stability}
\end{equation}
Equivalently, every claim in $\operatorname{Top}_k(P(Q))$ belongs to $H(Q)$.
\end{theorem}
\begin{proof}
Let $C=\operatorname{Top}_k(P(Q))$.
If $C\subseteq H(Q)$, every world retains $C$ and selects the same ranked context.
Conversely, take $c\in C\setminus H(Q)$.
Some world supports $c$, while another does not.
At most $k-1$ possible claims outrank $c$.
Thus every world supporting $c$ must select it.
The two worlds return different contexts.
For $k=0$, all contexts are empty.
\end{proof}

This test identifies uncertainty that affects the returned context.
Lower-ranked claims can remain uncertain without changing retrieval.
The test needs two filtered rankings and no enumeration of date assignments.
Its guarantee concerns retrieved evidence, rather than generated answers.

\subsection{A worked certificate}

Consider three claims with occurrence bounds $c=[0,2]$, $d=[1,4]$, and $e=[3,5]$.
Accept the directed replacements $d\to c$ and $e\to c$.
Claim $d$ can occur after $c$, but it need not do so.
Its earliest allowed start supplies $r_c=1$.
Claim $e$ must occur after $c$.
By time 5 it has occurred under every assignment, so it supplies $p_c=5$.
Table~\ref{tab:worked} shows the resulting decisions.

\begin{table}[t]
\centering\small
\begin{tabular}{@{}llll@{}}
\toprule
Query $Q$ & Possible & Guaranteed & Status\\
\midrule
$[0,0]$ & Yes & No & Unresolved\\
$[0,2]$ & Yes & Yes & Guaranteed\\
$[1,2]$ & Yes & No & Unresolved\\
$[2,4]$ & Yes & No & Unresolved\\
$[5,5]$ & No & No & Unsupported\\
\bottomrule
\end{tabular}
\caption{Illustrative decisions for $c=[0,2]$, $p_c=5$, and $r_c=1$. A larger query window can have guaranteed support even when its final point does not.}
\label{tab:worked}
\end{table}

The two witnesses preserve different decisions.
If we retain only $d$, time 5 incorrectly becomes possible for $c$.
If we retain only $e$, time 2 incorrectly becomes certain.
Thus neither one-witness subset reproduces all outputs.
This construction establishes the worst-case two-witness requirement.
The lower bound concerns retained source witnesses, rather than an unrestricted numerical encoding.

The example also separates date ambiguity from retrieval ambiguity.
For $Q=[0,2]$, rank $c$ before $d$ and set $k=1$.
Claim $c$ is guaranteed, so every timeline returns the same context.
Claim $d$ remains unresolved below the budget.
Increasing the budget to two makes the context depend on whether $d$ has started.
The stability test detects that change without enumerating timelines.

Now suppose later evidence narrows the same dates to $x_c=2$, $x_d=1$, and $x_e=5$.
Every refined date remains within its original bounds.
Claim $d$ can no longer replace $c$, since it occurs earlier.
Both certificates now identify $e$, and $c$ is active on $[2,5)$.
The window $[2,4]$ becomes guaranteed.
The point $[0,0]$ becomes unsupported.
The refinement removes possible support and resolves uncertainty into guaranteed support, as predicted by the corollary.

\subsection{Why the query tests are exact}

The possible-support boundary follows from a worst-case ordering test.
A witness whose earliest start exceeds $c$'s latest start must follow $c$.
Once that witness's latest date passes, it has retired $c$ in every timeline.
Before the first such boundary, each accepted witness can occur before $c$ or after the queried time.
Independent occurrence bounds permit these choices together.
This gives the possible point set $[\ell_c,p_c)$.

For guaranteed window support, first ensure that $c$ starts by the window's end in every timeline.
This is the condition $u_c\leq b$.
Any start inside the window immediately supplies support.
Failure therefore requires an earlier start followed by replacement before the window opens.
A witness can produce that failure exactly when $a>\ell_c$ and $r_c\leq a$.
Negating this condition gives the guaranteed-support test.

This argument explains why simply intersecting two intervals is insufficient.
Possible support has an ordinary pointwise interval.
Guaranteed window support instead depends on whether an entire timeline can avoid support throughout the query.
The compiler answers that question directly.

\subsection{Stable updates and inspectable decisions}

A certificate contains the target identifier, its date bounds, both minima, and their witness identifiers.
A retirement-based exclusion traces to its replacement witness.
A date-based exclusion traces to the target bounds.
Changing that decision updates only the target's two minima.
This is the dependency underlying Corollary~\ref{cor:uncertain-refinement}.
Transitive grouping creates a different dependency: one edge can merge entire histories.
Direct compilation avoids that propagation.

The error bound concerns eligibility and fixed-ranked contexts.
For $h$ changed directed pairs, at most $h$ target masks can change.
The selected sets differ by at most $2\min(k,h)$ items.
A removal can cause one new item to enter the context.
It can shift several positions within the returned list.
The bound concerns selected sets, rather than ordered positions.

Fix the existing date bounds and replacement decisions.
Adding claims whose lower date bounds exceed the query end preserves earlier support decisions.
Such a claim cannot start within that window or retire another claim by its end.
This follows from the same direct support definition.
If an added source instead reveals an earlier event, the historical mask can change.
The date bound determines the effect, rather than the ingestion order.

Finally, exact dates recover a familiar special case.
Possible and guaranteed support coincide.
One earliest direct replacement determines each claim's expiry.
The uncertain-date compiler extends this policy while preserving all permitted occurrence orders.
